% !TEX program = xelatex
% 海南热带海洋学院本科毕业论文（设计）模板
% 专业类别选项：science（理工，默认）/ humanities（文经管法）/ foreign（外语）/ art（艺术）
\ifdefined\HNTcategory\else\def\HNTcategory{science}\fi
\documentclass[\HNTcategory]{hntouthesis}

%====================== 论文信息（改这里） ======================
\renewcommand{\hnttitle}{单片机在测温技术中的应用}
\renewcommand{\hnttitleeng}{Design and Evaluation of a Temperature Measurement System}
\renewcommand{\hntauthor}{张\quad 三}
\renewcommand{\hntstudentid}{2022000001}
\renewcommand{\hntcollege}{海洋科学技术学院}
\renewcommand{\hntmajor}{电子信息工程}
\renewcommand{\hntclass}{2022 级 1 班}
\renewcommand{\hntadvisor}{李\quad 四}
\renewcommand{\hntadvisortitle}{副教授}
\renewcommand{\hntadvisorB}{}          % 第二位指导教师，没有就留空
\renewcommand{\hntadvisortitleB}{}
\renewcommand{\hntdate}{二〇二六年六月}
\renewcommand{\hntcohort}{26}          % 26 届
\ifhntforeign
  \renewcommand{\hnttitle}{外语研究论文写作与排版示例}
  \renewcommand{\hnttitleeng}{A Writing and Typesetting Example for a Foreign-Language Thesis}
  \renewcommand{\hntcollege}{外国语学院}
  \renewcommand{\hntmajor}{英语}
\fi
\ifhntart
  \renewcommand{\hnttitle}{毕业设计说明书与作品图版示例}
  \renewcommand{\hntcollege}{艺术学院（请核对）}
  \renewcommand{\hntmajor}{请填写专业方向}
\fi
\ifhnthuman
  \renewcommand{\hnttitle}{旅游服务研究论文写作与排版示例}
  \renewcommand{\hnttitleeng}{A Writing and Typesetting Example for Tourism Service Research}
  \renewcommand{\hntcollege}{请填写所在学院}
  \renewcommand{\hntmajor}{旅游管理（示例）}
\fi

\begin{document}
\ifhntforeign
  \linespread{1.52}\selectfont
\else
  \linespread{1.38}\selectfont
\fi
\setlength{\parindent}{2em}
\setlength{\parskip}{0pt}

%====================== 封面 ======================
\makehntcover
\makehntinnercover
\ifhntforeign\makehntengcover\fi
\hntbodylayout

\ifhntforeign
% 外语类单独使用英文正文，不将理工正文简单换字体。
\pagestyle{hntfront}
\begin{center}{\hei\zihao{-2}\hnttitle}\par\medskip{\hei\zihao{3}摘\quad 要}\end{center}
本文为外语类论文排版示例，不报告实证研究结果。模板展示中英文摘要、英文致谢、英文目录、英文正文及参考文献的组织顺序。正式写作应根据实际选题补充研究背景、理论框架、语料来源、分析方法、研究发现与局限，引用可核验的原始材料，不应保留本文件的示例性陈述。
\par\medskip\noindent{\hei\zihao{4}关键词：}论文写作，研究方法，格式示例
\clearpage
\begin{center}{\rmfamily\bfseries\zihao{-2}\hnttitleeng}\par\medskip
Author: \hntauthoreng\par Supervisor: \hntadvisoreng\par\medskip
{\bfseries\zihao{3}ABSTRACT}\end{center}
This document demonstrates the organization of a foreign-language thesis. It includes bilingual abstracts, acknowledgements, a table of contents, a method section, and a reference list. It does not present an empirical study. All examples must be replaced with the author's own research, supported by verifiable sources and a clearly described method.
\par\medskip\noindent\textbf{Keywords:} thesis writing, research methods, typesetting
\clearpage\section*{Acknowledgements}
This page precedes the contents in the supplied official foreign-language example. Replace this text with a personal and accurate acknowledgement of the people and institutions that supported the work.
\clearpage
\renewcommand{\contentsname}{CONTENTS}
\renewcommand{\cfttoctitlefont}{\hfill\rmfamily\bfseries\zihao{2}}
\renewcommand{\cftsecfont}{\rmfamily\bfseries\zihao{4}}
\renewcommand{\cftsubsecfont}{\rmfamily\zihao{-4}}
\renewcommand{\cftsubsubsecfont}{\rmfamily\zihao{-4}}
\renewcommand{\cftdotsep}{4.5}
\tableofcontents
\clearpage\pagenumbering{arabic}\pagestyle{hntmain}
\section{Introduction}
\subsection{Research question}
A thesis should define a specific question, explain why it matters, and describe the evidence needed to answer it. The present document is a formatting example only. It makes no claim about language use, translation quality, or cultural differences.
\subsection{Scope and contribution}
Define the language pair, genre, period, and sampling frame before collecting material. A contribution may consist of a transparent corpus, an explicit analytical procedure, or a carefully supported interpretation. Avoid presenting stereotypes as established cultural facts.
\section{Literature Review}
\subsection{Concepts and evidence}
Organize the review around competing explanations and methodological choices. Distinguish the claims made by a source from the evidence it provides. Real publications are included here solely to demonstrate citation commands\cite{knuth1984,lamport1994}; they are not a literature review of translation studies.
\section{Research Design}
\subsection{Materials and sampling}
Report how texts were selected, excluded, and archived. If human participants are involved, describe the relevant ethics approval and consent procedures. Do not infer that publicly visible personal data may be reused without restrictions.
\subsubsection{Coding procedure}
Define categories before interpreting examples, document ambiguous cases, and explain how disagreements were resolved.
\paragraph{A fourth-level heading}
The supplied Word example uses four levels of decimal headings; only three levels are listed in the contents.
\section{Analysis and Discussion}
Link each analytical claim to specific examples and consider alternative explanations. A table can summarize coding categories, but it cannot replace a reasoned interpretation. No empirical findings are supplied in this template.
\section{Conclusion}
Answer the research question with appropriately bounded claims, explain limitations, and identify realistic future work. Do not introduce evidence that has not been discussed in the body.
\clearpage\begingroup
\renewcommand{\bibsection}{\section*{References}\addcontentsline{toc}{section}{References}}
\renewcommand{\bibfont}{\rmfamily\zihao{5}}
\bibliographystyle{gbt7714-2015-numeric}\bibliography{refs}
\endgroup

\else
\ifhntart
%====================== 艺术类：设计说明 ======================
\newpage
\pagenumbering{arabic}
\pagestyle{hntmain}
\section*{设计说明}
\addcontentsline{toc}{section}{设计说明}
\section{设计背景与思路}
\subsection{设计背景}
此处撰写毕业设计作品的创作背景、选题依据与设计目标。

\subsection{设计思路}
此处说明作品的主题、材料、工艺、空间或视觉表达方案，可配合图片说明。

\section{作品展示与分析}
此处放置作品照片或效果图，并从设计理念、形式语言、实施过程等方面进行分析。
\section{设计过程与实施}
\subsection{资料调研与方案演化}
记录素材来源、授权范围及设计约束，用草图、迭代记录和最终作品说明设计决策。官方艺术类文件明确允许说明书版式自定，因此本节是可修改的设计报告结构，不是学校规定的章节清单。
\subsection{作品图版示例}
以下为矢量构图占位示例，不是毕业设计作品；请替换为具有使用权的作品照片或效果图。
\begin{figure}[htbp]\centering
\begin{tikzpicture}
\fill[black!8] (0,0) rectangle (9,4);
\draw[very thick] (0.6,0.6) rectangle (4,3.4);
\draw[very thick] (5,0.6)--(8.4,3.4) (5,3.4)--(8.4,0.6);
\node at (4.5,-0.4) {图版占位：替换为实际作品};
\end{tikzpicture}
\caption{构图占位示例（非作品成果）}
\end{figure}
\section{反思与成果归档}
说明设计目标的达成情况、制作成本、材料限制、改进方向和作品版权。按官方附件要求另备 A4 彩色作品图版；研究涉及非遗图案、人物肖像或第三方素材时，应注明出处并核对授权。


\newpage
\section*{致\quad 谢}
\addcontentsline{toc}{section}{致谢}
感谢指导教师和同学在毕业设计过程中给予的帮助。

\vspace{2cm}
\begin{flushright}
学生签名：\underline{\makebox[3cm]{}}\\[1.2em]
日\quad 期：\underline{\makebox[3cm]{}}\ 年\ \underline{\makebox[1cm]{}}\ 月\ \underline{\makebox[1cm]{}}\ 日
\end{flushright}
\else
%====================== 中文摘要 ======================
\newpage
\pagestyle{hntfront}
\begin{center}
{\hei\zihao{-2} \hnttitle}
\end{center}
\vspace{0.4em}
{\song\zihao{-4}\hfill 作者：\hntauthor\par}
{\song\zihao{-4}\hfill 导师：\hntadvisor\par}
\vspace{0.8em}
\begin{center}
{\hei\zihao{-3} 摘\quad 要}
\end{center}
\vspace{0.4em}

{\song\zihao{-4}
\ifhnthuman
本文件为旅游服务研究论文的写作与排版示例，不报告真实问卷、访谈或企业调查。模板以服务信息传递为线索，展示问题定义、文献综述、抽样设计、观察维度、资料编码、证据报告和讨论的组织方式，并提供可编辑的三线表、流程图、公式与每页重新编号脚注。文中没有预填受访者数量、满意度比例或统计显著性；正式论文应根据真实且获准使用的资料独立完成，准确区分参与者陈述、观察事实和作者解释，并说明抽样范围、隐私保护与研究局限。摘要应概括实际研究的主要内容和结果，不写成目录式介绍，也不将一般建议当作已经证实的结论。官方文经管法类样例建议摘要约250至350字，关键词3至5个；最终仍需核对所在学院当年通知。
\else
本文件为毕业论文排版与写作示例，不报告真实实验成果。模板以温度采集系统为线索，展示研究问题、技术路线、实验设计、误差分析和结论的组织方法，并提供可直接修改的结构图、曲线图、三线表、数学公式、脚注及程序清单。正文中的所有测量数据均为人为构造的排版示例，不能作为科研证据。正式论文应以实际开展的研究为基础，摘要按照目的、方法、结果和结论概括主要内容，补入真实样本量、测试条件和关键结果，避免使用没有数据支持的“高精度”“显著提升”等判断。理工及文经管法类官方样例建议摘要约250至350字，并列出3至5个关键词；具体提交要求以学院当年通知为准。
\fi
\par}

\vspace{0.6em}
\noindent{\hei\zihao{-4} 关键词：}{\hei\zihao{-4}\ifhnthuman 旅游服务；研究设计；资料分析；排版示例\else 单片机；温度测量；DS18B20；液晶显示\fi}

%====================== 英文摘要 ======================
\newpage
\begin{center}
{\rmfamily\zihao{-2}\bfseries \hnttitleeng}
\end{center}
\vspace{0.4em}
{\rmfamily\zihao{-4}\hfill Author: \hntauthoreng\par}
{\rmfamily\zihao{-4}\hfill Tutor: \hntadvisoreng\par}
\vspace{0.8em}
\begin{center}
{\rmfamily\zihao{-3}\bfseries ABSTRACT}
\end{center}
\vspace{0.4em}

{\rmfamily\zihao{-4}
\ifhnthuman
This document demonstrates the structure and typesetting of a tourism service research thesis. It does not report a survey, an interview study, or a company investigation. The examples cover problem definition, sampling, coding, evidence reporting, and discussion, together with tables, diagrams, equations, and footnotes. No participant counts, satisfaction rates, or statistical findings are supplied. Replace all examples with independently conducted research and accurately distinguish observations, participant statements, and the author's interpretations. The final abstract must summarize the actual study rather than present a list of chapters.
\else
This document is a typesetting and writing example rather than a report of experimental findings. A temperature measurement system is used to demonstrate research questions, system architecture, experimental design, and error analysis. The examples cover diagrams, plots, tables, equations, footnotes, and source code. All numerical results are synthetic and must be replaced with independently obtained research data. The final abstract should accurately describe the purpose, methods, results, and conclusions of the actual study without unsupported claims of accuracy or improvement.
\fi
\par}

\vspace{0.6em}
\noindent{\rmfamily\zihao{-4}\bfseries Keywords:} {\rmfamily\zihao{-4}\ifhnthuman tourism service; research design; data analysis; typesetting\else microcontroller; temperature measurement; DS18B20; LCD\fi}

%====================== 目录 ======================
\newpage
\pagestyle{hntfront}
\tableofcontents
\thispagestyle{hntfront}

%====================== 正文 ======================
\newpage
\pagenumbering{arabic}
\pagestyle{hntmain}

\ifhnthuman\section{绪论}
\subsection{研究问题与研究边界}
旅游服务研究宜从具体场景提出问题，明确服务环节、参与者和分析单位。本例以服务信息传递为写作线索，演示文经管法类编号和研究论文结构，不报告真实问卷、访谈或企业资料。城市、机构和调查人数应在研究完成后如实填写。
\subsection{文献综述与概念界定}
综述需要比较概念、方法、证据及局限，不能把一般经验直接称为理论。资料应追溯至原始出版物，记录版本、页码和允许用途。本例使用真实出版物演示引文排版\cite{knuth1984,lamport1994}，不把它们冒充为旅游服务研究综述。
\subsubsection{研究对象与分析单位}
明确分析对象是游客个体、单次服务事件、团队还是机构。同一游客多次反馈不自动等于多个独立样本；便利样本也不代表所有游客。涉及个人资料时，应取得适当许可并匿名化\footnote{本例没有接触真实参与者。脚注只演示同页补充说明与每页重置编号。}。
\paragraph{四级标题示例}
本类编号采用“一、㈠、1、⑴”，目录仅列到三级。不要因模板提供四级编号就把每个自然段都设为标题。
\section{研究设计}
\subsection{资料来源与抽样}
报告资料来源、时间范围、选择方式和排除规则。访谈提纲应经过试访修订，问卷指标应说明来源及适用条件，不能随意拼接后声称具有经过验证的信效度。
\subsection{分析维度与编码规则}
表~\ref{tab:tourism} 是观察维度表的结构示例，并非实际调查结论。正式研究需给出类别边界、难以判断的情况及复核过程。
\begin{table}[htbp]\centering\caption{服务信息研究维度示例}\label{tab:tourism}
\begin{tabularx}{\linewidth}{lXX}\toprule 维度 & 记录内容 & 解释边界\\\midrule
信息获取 & 来源、时间点与可访问性 & 区分公开材料与个人体验\\
服务交互 & 具体事件与沟通过程 & 不以单例推断总体水平\\
问题处理 & 反馈渠道与响应过程 & 记录上下文与反例\\\bottomrule\end{tabularx}
\end{table}
\subsection{研究流程}
图~\ref{fig:tourism} 展示从问题定义到结果解释的组织方式。每一步均应留下过程材料，以便核查。
\begin{figure}[htbp]\centering
\begin{tikzpicture}[node distance=6mm,b/.style={draw,minimum width=65mm,minimum height=9mm,font=\small},>=Stealth]
\node[b](a){明确问题与取得必要许可};\node[b,below=of a](b){试用工具与收集资料};
\node[b,below=of b](c){整理、匿名化与编码复核};\node[b,below=of c](d){解释证据与讨论局限};
\draw[->](a)--(b);\draw[->](b)--(c);\draw[->](c)--(d);
\end{tikzpicture}\caption{研究流程示例}\label{fig:tourism}
\end{figure}
\subsection{统计口径与质性分析}
采用比例时，应说明分母、缺失与多选情况；采用质性分析时，应呈现主题形成、反例和解释路径。若某类事件数为 $n_k$，满足纳入条件的总事件数为 $N$，其比例示例如下：
\begin{equation}p_k=\frac{n_k}{N}\times100\%.\end{equation}
没有实际数据时不填入百分比，更不能虚构统计显著性。统计指标只是工具，不能替代研究问题和语境分析。
\section{结果报告}
\subsection{描述性结果}
正式报告应先说明资料数量、覆盖范围及排除情况，再按研究问题呈现发现。摘录应匿名并保留内部编号，不应断章取义。本示例没有真实调查，因此只提供报告结构，没有“发现的问题”可供引用。
\subsection{证据、观点与解释}
区分研究者直接观察、参与者陈述和研究者解释。负面反馈不等于所有服务均有问题，满意评价也不能证明某项管理措施产生因果效果。图表与文字应互相补充，不重复堆放同一信息。
\section{讨论与建议}
\subsection{与已有研究的关系}
比较研究结果与既有研究的一致或差异，说明场景、抽样和测量方法可能带来的影响。若提出服务改进建议，应明确实施主体、成本、风险及可检验目标，避免空泛对策。
\subsection{局限与后续工作}
说明样本选择、研究者参与、资料缺失和隐私限制等问题。未来计划不能写成已完成的成果。本工作稿不支持任何关于旅游服务质量的实证判断。
\else% 本文是写作与排版操作示例，不是可提交的研究论文。
\section{绪论}
\subsection{研究问题与研究边界}
论文应首先说明需要解决的具体问题，而不是直接罗列器件或软件名称。本例的问题是：在给定测温范围、采样周期和使用环境下，如何组织数据采集、异常检测与结果显示，并用可追溯的方法评价测量误差。研究边界应同时给出，包括温度范围、供电条件、参考仪器以及不在本研究范围内的场景。

资料检索应优先使用器件数据手册、标准和原始研究论文。引用前应核对作者、题名、版本与可访问地址，不应将模板中的占位文献作为已经阅读的材料。本例使用真实出版的教材和技术报告演示文献排版\cite{knuth1984,lamport1994}。
\subsection{文献综述的组织}
综述宜围绕研究问题比较已有方法的假设、证据及限制，而不只是按作者排列摘要。可建立“研究对象—方法—数据—结论—局限”的文献矩阵；同类结果的测试条件不一致时，不宜直接以数值高低判断优劣。文献管理与正文写作应同步进行，记录实际使用的页码及访问日期。
\subsubsection{直接引用与解释性脚注}
参考文献用于标识已有工作的出处，解释性脚注用于补充不适合放在正文中的信息，两者不能混为一谈。这里演示同页重新编号的脚注\footnote{本脚注用于说明排版方法，不包含研究证据。脚注应随其引用位置出现在同一页。}。连续引用可使用 \verb|\cite{knuth1984,lamport1994}|，交叉引用应使用标签而非手写页码。
\paragraph{四级标题示例}
理工类使用四级阿拉伯数字编号，文经管法类使用“一、㈠、1、⑴”的层级形式。四级标题不进入三级目录；正式写作应避免为了展示编号而过度切分内容。
\subsection{研究内容与技术路线}
研究任务可分解为需求界定、方案比较、系统实现、试验设计及结果讨论。每个任务应有相应产物，例如设计指标、方案选择依据、接口定义、原始记录和可复核的计算过程。下文展示如何在同一文档中引用图、表和公式，不用外部图片文件也能编译。

\section{总体方案与方法}
\subsection{系统结构与接口}
图~\ref{fig:architecture} 展示数据从采集端到记录端的流向。正式系统图应标出关键接口、采样频率和异常处理路径，避免只画模块而不说明模块间传递什么信息。
\begin{figure}[htbp]
\centering
\begin{tikzpicture}[node distance=5mm and 5mm,box/.style={draw,rounded corners,minimum width=23mm,minimum height=11mm,align=center,font=\small},>=Stealth]
\node[box] (sensor) {温度采集};
\node[box,right=of sensor] (check) {数据校验};
\node[box,right=of check] (process) {换算与滤波};
\node[box,below=of process] (record) {显示与记录};
\draw[->] (sensor)--(check);\draw[->] (check)--(process);\draw[->] (process)--(record);
\end{tikzpicture}
\caption{数据处理结构示意图（方法示例）}\label{fig:architecture}
\end{figure}
\subsection{方案比较与选择依据}
方案比较应使用与需求直接相关的指标，而不是给候选方案任意打分。对无法量化的约束可以说明判定条件和证据来源。表~\ref{tab:requirements} 为需求追踪表的结构示例，不代表某项器件的实际性能。
\begin{table}[htbp]
\centering\caption{需求与验证方法对应表}\label{tab:requirements}
\small
\begin{tabularx}{\linewidth}{lXX}
\toprule 需求项 & 说明 & 验证方法\\\midrule
测量范围 & 在任务书中明确上下限 & 分温度点测试并保存原始记录\\
采样周期 & 给出采集与显示更新的区别 & 时间戳统计与逻辑分析\\
异常处理 & 识别断连、超时及无效读数 & 主动故障注入与恢复测试\\
重复性 & 相同条件下的离散程度 & 重复测量与标准差估计\\\bottomrule
\end{tabularx}
\end{table}
\subsection{误差定义与数学表达}
设第 $i$ 次测量值为 $y_i$，相应参考值为 $x_i$，误差定义为 $e_i=y_i-x_i$。均方根误差的计算示例如式~\eqref{eq:rmse} 所示：
\begin{equation}\label{eq:rmse}
\operatorname{RMSE}=\sqrt{\frac{1}{n}\sum_{i=1}^{n}(y_i-x_i)^2}.
\end{equation}
式中，$n$ 为有效测量次数。误差指标本身不能代替测量不确定度分析；参考仪器、环境波动和重复测量的相关性也应说明。国际计量联合委员会的不确定度指南可作为相关方法的阅读材料\cite{jcgm2008}。

\section{实现与可复现性}
\subsection{模块实现与异常处理}
实现章节应描述输入、输出、约束及错误分支，而不是把完整代码粘贴成正文。采集模块至少需要说明通信失败时如何返回状态，数据校验模块说明如何区分无效数据与真实极值，显示模块说明异常状态的呈现方式。关键接口的选择应能够追溯至前文需求。
\subsubsection{参数配置}
采样周期、平滑窗口和报警阈值应作为显式参数记录。改变参数进行比较时，应保持其他条件不变；如果多项条件同时改变，需要讨论混杂因素。论文应给出软件版本、依赖版本、运行命令和原始数据存放位置，保证读者可以重建分析过程。
\subsection{最小计算示例}
下面的程序只演示误差计算，不包含真实采集驱动；输入数组明确为合成数据。正式论文应将大段代码放在附录，并在正文解释其作用和验证方法。
\begin{lstlisting}[language=Python,caption={误差计算示例（输入为合成数据）}]
from math import sqrt

reference = [20.0, 25.0, 30.0, 35.0]
observed = [20.1, 24.9, 30.2, 35.0]
errors = [y - x for x, y in zip(reference, observed)]
rmse = sqrt(sum(e * e for e in errors) / len(errors))
print(f"RMSE = {rmse:.3f}")
\end{lstlisting}

\section{试验设计与结果报告}
\subsection{试验条件与记录规范}
真实试验应明确仪器型号、校准情况、测试环境、稳定判据和每个条件下的重复次数。记录中应保留失败数据及排除理由，不能只保留看起来较好的结果。研究涉及人体、访谈或敏感信息时，还应取得相应审批和知情同意，不能用“模板示例”替代必要程序。
\subsection{合成数据的表格与曲线示例}
表~\ref{tab:synthetic} 和图~\ref{fig:errors} 仅供观察版式；没有开展实验，也不支持任何性能结论。表中单位写入表头，数据保留的小数位应与实际测量精度一致。
\begin{table}[htbp]
\centering\caption{合成测量数据（非实验结果）}\label{tab:synthetic}
\begin{tabular}{rrrr}\toprule 序号 & 参考值/℃ & 示例读数/℃ & 误差/℃\\\midrule
1 &20.0&20.1&0.1\\2 &25.0&24.9&-0.1\\3 &30.0&30.2&0.2\\4 &35.0&35.0&0.0\\\bottomrule
\end{tabular}
\end{table}
\begin{figure}[htbp]\centering
\begin{tikzpicture}
\begin{axis}[width=0.9\linewidth,height=52mm,xlabel={参考温度/℃},ylabel={误差/℃},xmin=19,xmax=36,ymin=-0.15,ymax=0.25,grid=major,tick label style={font=\small},label style={font=\small}]
\addplot[black,mark=*] coordinates {(20,0.1)(25,-0.1)(30,0.2)(35,0)};
\end{axis}
\end{tikzpicture}
\caption{合成数据误差曲线（仅演示坐标与单位）}\label{fig:errors}
\end{figure}
\subsection{讨论与局限}
讨论应解释结果与研究问题之间的关系，区分观察事实、计算结果和作者推断。对于样本量不足、参考仪器精度有限或条件覆盖不全等问题，应说明其可能影响，不宜只写“仍有不足”。由于本模板没有真实研究，不能讨论统计显著性或声称方案优于已有工作。
\subsection{提交前复核}
检查目录能否正确跳转，图表与公式是否有正文引导，是否存在未定义引用，附录是否与正文呼应。最后核对学院当年模板、签字要求、查重文件格式及归档材料；PDF 编译成功不意味着学校接受 PDF 替代 Word。
\fi

\newpage
\section*{结\quad 论}
\addcontentsline{toc}{section}{结论}
本示例完成了论文主要排版元素的串联，展示了从问题定义、方案描述到结果报告的组织方式。上述示例没有实施真实实验或调查，不能得出技术性能或服务质量结论。正式论文应逐项回答研究问题，说明证据支持的结论、适用条件及尚未解决的局限，不应把未来计划写成已经完成的成果。

%====================== 参考文献 ======================
\newpage
\begingroup
\renewcommand{\bibsection}{%
  \phantomsection
  \addcontentsline{toc}{section}{参考文献}%
  \begin{center}\hei\zihao{3}参考文献\end{center}
  \vspace{0.5\baselineskip}%
}
\renewcommand{\bibfont}{\song\zihao{5}}
\bibliographystyle{gbt7714-2015-numeric}
\bibliography{refs}
\endgroup

%====================== 附录 ======================
\newpage
\section*{附\quad 录}
\addcontentsline{toc}{section}{附录}
{\song\zihao{5}\linespread{1.0}\selectfont
此处放置需要收录于论文中、但不适合写入正文的附加数据、程序清单、详细公式推导等内容。
\par\medskip\noindent 附录1\quad 数据字典示例\par
\begin{longtable}{lll}
\toprule 字段 & 含义 & 单位或取值\\\midrule\endfirsthead
\multicolumn{3}{c}{附录1（续）}\\\toprule 字段 & 含义 & 单位或取值\\\midrule\endhead
sample\_id & 样本标识 & 不重复编号\\
reference & 参考温度 & ℃\\
observed & 采集读数 & ℃\\
status & 采集状态 & valid / timeout / invalid\\
timestamp & 采样时间 & ISO 8601 时间戳\\\bottomrule
\end{longtable}
附录宜提供可复核的补充材料，例如原始字段定义、访谈提纲、详细推导和运行说明。长表使用 longtable 环境，续页会重复表头；不要通过截图插入无法搜索的表格。

}

%====================== 致谢 ======================
\newpage
\section*{致\quad 谢}
\addcontentsline{toc}{section}{致谢}
感谢指导教师在选题、方案设计和论文写作过程中给予的悉心指导，感谢同学和家人在学习期间的支持与帮助。

\vspace{2cm}
\begin{flushright}
学生签名：\underline{\makebox[3cm]{}}\\[1.2em]
日\quad 期：\underline{\makebox[3cm]{}}\ 年\ \underline{\makebox[1cm]{}}\ 月\ \underline{\makebox[1cm]{}}\ 日
\end{flushright}
\fi
\fi

\end{document}
